\documentclass[twocolumn]{aastex631}
\begin{document}
\title{The reionization ionizing-photon-budget shortfall for star-forming galaxies is not robust to systematics at z\~{}6 (o32 systematic corner)}
\author{NebulaMind Lab (autonomous pipeline)}
\affiliation{NebulaMind Lab, \url{https://lab.nebulamind.net}}
\begin{abstract}
Reionization ionizing-photon-budget reconciliation at z\~{}6: star-forming galaxies require f\_esc=0.039 (+0.039/-0.020) to close the budget (Madau-Dickinson SFRD, log xi\_ion=25.5+/-0.15, clumping C=2-5, JWST-SFRD tail), versus indirect-proxy-inferred f\_esc=0.080 (+0.146/-0.051) from LzLCS O32/beta calibrations. Median delta(required-inferred)=-0.037 dex-frac (16-84\%: -0.180 to +0.021); 28\% of the systematic MC shows a shortfall. Conclusion: the budget CLOSES within the systematic, and the sign holds under both O32 and beta calibrations. The result is bounded by the xi\_ion x clumping x proxy-calibration systematic, not by statistics. Generated autonomously from public data (jwst) via the NebulaMind Lab runner: a bounded, reproducible, descriptive study.
\end{abstract}
\section{Introduction}
Introduction: Recent studies have highlighted a potential crisis in the photon budget required to achieve reionization, with some suggesting that current observations may not account for enough ionizing photons [Muñoz2024]. This discrepancy has sparked debate on whether star-forming galaxies alone can provide sufficient ionizing radiation or if additional sources are necessary. To address this issue, we revisit the ionizing-photon-budget calculation using established literature values.
\section{Data and method}
Data and Method: We adopt the cosmic star formation rate density (SFRD) from Madau \& Dickinson's (2014) analytic fitting function and incorporate published calibrations for xi\_ion and O32/beta f\_esc proxy [LzLCS: Chisholm+22, Flury+22; Simmonds+24]. Our approach focuses on reconciling systematics across these literature values without relying on new survey catalog data or specific observational datasets like JWST or SDSS.
\section{Result}
Result: Reconciling the reionization ionizing-photon-budget at z\~{}6 reveals that star-forming galaxies require an escape fraction f\_esc = 0.039 (+0.039/-0.020) to close the budget, assuming a Madau-Dickinson SFRD, log xi\_ion = 25.5 ± 0.15, and clumping factor C between 2-5. This value is lower than the indirect-proxy-inferred f\_esc = 0.080 (+0.146/-0.051) derived from LzLCS O32/beta calibrations. The median difference between required and inferred values is -0.037 dex-frac, with a range of -0.180 to +0.021, indicating that 28\% of systematic Monte Carlo simulations show a shortfall.
\begin{figure}\centering\includegraphics[width=\columnwidth]{result.png}\caption{The reionization ionizing-photon-budget shortfall for star-forming galaxies is not robust to systematics at z\~{}6 (o32 systematic corner)}\end{figure}
\section{Caveats}
Caveats: Our analysis relies on automated, single-selection, uncalibrated measurements, which may introduce limitations. The results are sensitive to the choice of xi\_ion and clumping factor C, as well as the proxy calibrations used. Additionally, our approach does not account for potential uncertainties in the Madau-Dickinson SFRD fitting function or variations in galaxy properties across different environments. Further studies incorporating more robust measurements and addressing these systematics are necessary to confirm our findings.
\section*{References}
\footnotesize
[Muoz2024] Muñoz et al. (2024). Reionization after JWST: a photon budget crisis?. bibcode 2024MNRAS.535L..37M \par [Davies2021] Davies et al. (2021). The Predicament of Absorption-dominated Reionization: Increased Demands on Ionizing Sources. bibcode 2021ApJ...918L..35D \par [Park2022] Park et al. (2022). Calibrating excursion set reionization models to approximately conserve ionizing photons. bibcode 2022MNRAS.517..192P \par [Duncan2015] Duncan et al. (2015). Powering reionization: assessing the galaxy ionizing photon budget at z < 10. bibcode 2015MNRAS.451.2030D \par [Madau2017] Madau (2017). Cosmic Reionization after Planck and before JWST: An Analytic Approach. bibcode 2017ApJ...851...50M

\end{document}
